\title{2025 年高考数学 全国II卷}
\date{}
\maketitle

\section*{一、单项选择题}
本题共8小题，每小题5分，共40分.在每小题给出的选项中，只有一项是符合题目要求的.

\begin{question}[points=5]
样本数据 $2$，$8$，$14$，$16$，$20$ 的平均数为（ ）
\begin{choices}
  \item $8$
  \item $9$
  \item $12$
  \item $18$
\end{choices}
\begin{solution}
\textbf{答案：}$C$

\textbf{分析：}由平均数的计算公式即可求解.

\textbf{详解：}样本数据 $2,8,14,16,20$ 的平均数为 $\frac{2+8+14+16+20}{5}=\frac{60}{5}=12$.
\end{solution}
\end{question}

\begin{question}[points=5]
已知 $z=1+\mathrm{i}$，则 $\frac{1}{z-1}=$（ ）
\begin{choices}
  \item $-\mathrm{i}$
  \item $\mathrm{i}$
  \item $-1$
  \item $1$
\end{choices}
\begin{solution}
\textbf{答案：}$A$

\textbf{分析：}由复数除法即可求解.

\textbf{详解：}因为 $z=1+\mathrm{i}$，所以 $\frac{1}{z-1}=\frac{1}{1+\mathrm{i}-1}=\frac{1}{\mathrm{i}}=-\mathrm{i}$.
\end{solution}
\end{question}

\begin{question}[points=5]
已知集合 $A=\{-4,0,1,2,8\}$，$B=\{x\mid x^3=x\}$，则 $A\cap B=$（ ）
\begin{choices}
  \item $\{0,1,2\}$
  \item $\{1,2,8\}$
  \item $\{2,8\}$
  \item $\{0,1\}$
\end{choices}
\begin{solution}
\textbf{答案：}$D$

\textbf{分析：}求出集合 $B$ 后结合交集的定义可求 $A\cap B$.

\textbf{详解：}$B=\{x\mid x^3=x\}=\{0,-1,1\}$，故 $A\cap B=\{0,1\}$.
\end{solution}
\end{question}

\begin{question}[points=5]
不等式 $\frac{x-4}{x-1}\ge 2$ 的解集是（ ）
\begin{choices}
  \item $\{x\mid -2\le x\le 1\}$
  \item $\{x\mid x\le -2\}$
  \item $\{x\mid -2\le x<1\}$
  \item $\{x\mid x>1\}$
\end{choices}
\begin{solution}
\textbf{答案：}$C$

\textbf{分析：}移项后转化为求一元二次不等式的解即可.

\textbf{详解：}$\frac{x-4}{x-1}\ge 2$ 即为 $\frac{x+2}{x-1}\le 0$，即 $\begin{cases}(x+2)(x-1)\le 0\\x-1\neq 0\end{cases}$，故 $-2\le x<1$，解集为 $[-2,1)$.
\end{solution}
\end{question}

\begin{question}[points=5]
在 $\triangle ABC$ 中，$BC=2$，$AC=1+\sqrt{3}$，$AB=\sqrt{6}$，则 $A=$（ ）
\begin{choices}
  \item $45^\circ$
  \item $60^\circ$
  \item $120^\circ$
  \item $135^\circ$
\end{choices}
\begin{solution}
\textbf{答案：}$A$

\textbf{分析：}由余弦定理 $\cos A=\frac{AB^2+AC^2-BC^2}{2AB\cdot AC}$ 直接计算求解即可.

\textbf{详解：}$\cos A=\frac{(\sqrt{6})^2+(1+\sqrt{3})^2-2^2}{2\cdot\sqrt{6}\cdot(1+\sqrt{3})}=\frac{6+1+2\sqrt{3}+3-4}{2\sqrt{6}(1+\sqrt{3})}=\frac{6+2\sqrt{3}}{2\sqrt{6}(1+\sqrt{3})}=\frac{\sqrt{2}}{2}$，又 $0^\circ<A<180^\circ$，所以 $A=45^\circ$.
\end{solution}
\end{question}

\begin{question}[points=5]
设抛物线 $C:y^2=2px(p>0)$ 的焦点为 $F$，点 $A$ 在 $C$ 上，过 $A$ 作 $C$ 的准线的垂线，垂足为 $B$，若直线 $BF$ 的方程为 $y=-2x+2$，则 $|AF|=$（ ）
\begin{choices}
  \item $3$
  \item $4$
  \item $5$
  \item $6$
\end{choices}
\begin{solution}
\textbf{答案：}$C$

\textbf{分析：}先由直线 $l_{BF}$ 求出焦点 $F$ 和 $p$ 即抛物线 $C$ 的方程，进而依次得抛物线的准线方程和点 $B$，从而可依次求出 $y_A$ 和 $x_A$，再由焦半径公式即可得解.

\textbf{详解：}对 $l_{BF}:y=-2x+2$，令 $y=0$，则 $x=1$，所以 $F(1,0)$，$p=2$ 即抛物线 $C:y^2=4x$，故抛物线的准线方程为 $x=-1$，故 $B(-1,4)$，则 $y_A=4$，代入抛物线 $C:y^2=4x$ 得 $x_A=4$，所以 $AF=AB=x_A+\frac{p}{2}=4+1=5$.
\end{solution}
\end{question}

\begin{question}[points=5]
记 $S_n$ 为等差数列 $\{a_n\}$ 的前 $n$ 项和，若 $S_3=6$，$S_5=-5$，则 $S_6=$（ ）
\begin{choices}
  \item $-20$
  \item $-15$
  \item $-10$
  \item $-5$
\end{choices}
\begin{solution}
\textbf{答案：}$B$

\textbf{分析：}由等差数列前 $n$ 项和公式结合题意列出关于首项 $a_1$ 和公差 $d$ 的方程求出首项 $a_1$ 和公差 $d$，再由等差数列前 $n$ 项和公式即可计算求解.

\textbf{详解：}设等差数列 $\{a_n\}$ 的公差为 $d$，则由题可得 $\begin{cases}3a_1+3d=6\\5a_1+10d=-5\end{cases}\Rightarrow\begin{cases}d=-3\\a_1=5\end{cases}$，所以 $S_6=6a_1+15d=6\times5+15\times(-3)=-15$.
\end{solution}
\end{question}

\begin{question}[points=5]
已知 $0<\alpha<\pi$，$\cos\frac{\alpha}{2}=\frac{\sqrt{5}}{5}$，则 $\sin\left(\alpha-\frac{\pi}{4}\right)=$（ ）
\begin{choices}
  \item $\frac{\sqrt{2}}{10}$
  \item $\frac{\sqrt{2}}{5}$
  \item $\frac{3\sqrt{2}}{10}$
  \item $\frac{7\sqrt{2}}{10}$
\end{choices}
\begin{solution}
\textbf{答案：}$D$

\textbf{分析：}利用二倍角余弦公式得 $\cos\alpha=-\frac{3}{5}$，则 $\sin\alpha=\frac{4}{5}$，最后再根据两角差的正弦公式即可得到答案.

\textbf{详解：}$\cos\alpha=2\cos^2\frac{\alpha}{2}-1=2\times\left(\frac{\sqrt{5}}{5}\right)^2-1=-\frac{3}{5}$，因为 $0<\alpha<\pi$，则 $\frac{\pi}{2}<\alpha<\pi$，则 $\sin\alpha=\sqrt{1-\cos^2\alpha}=\sqrt{1-\left(-\frac{3}{5}\right)^2}=\frac{4}{5}$，则 $\sin\left(\alpha-\frac{\pi}{4}\right)=\sin\alpha\cos\frac{\pi}{4}-\cos\alpha\sin\frac{\pi}{4}=\frac{4}{5}\times\frac{\sqrt{2}}{2}-\left(-\frac{3}{5}\right)\times\frac{\sqrt{2}}{2}=\frac{7\sqrt{2}}{10}$.
\end{solution}
\end{question}

\section*{二、多项选择题}
本题共3小题，共18分.在每小题给出的选项中，有多项符合题目要求.

\begin{question}[points=6]
记 $S_n$ 为等比数列 $\{a_n\}$ 的前 $n$ 项和，$q$ 为 $\{a_n\}$ 的公比，$q>0$，若 $S_3=7$，$a_3=1$，则（ ）
\begin{choices}
  \item $q=\frac{1}{2}$
  \item $a_5=\frac{1}{9}$
  \item $S_5=8$
  \item $a_n+S_n=8$
\end{choices}
\begin{solution}
\textbf{答案：}$AD$

\textbf{分析：}对A，根据等比数列通项公式和前 $n$ 项和公式得到方程组，解出 $a_1,q$，再利用其通项公式和前 $n$ 项和公式一一计算分析即可.

\textbf{详解：}对A，由题意得 $\begin{cases}a_1q^2=1\\a_1+a_1q+a_1q^2=7\end{cases}$，结合 $q>0$，解得 $\begin{cases}a_1=4\\q=\frac{1}{2}\end{cases}$ 或 $\begin{cases}a_1=9\\q=-\frac{1}{3}\end{cases}$（舍去），故A正确；对B，则 $a_5=a_1q^4=4\times\left(\frac{1}{2}\right)^4=\frac{1}{4}$，故B错误；对C，$S_5=\frac{a_1(1-q^5)}{1-q}=\frac{4\times\left(1-\frac{1}{32}\right)}{1-\frac{1}{2}}=\frac{31}{4}$，故C错误；对D，$a_n=4\times\left(\frac{1}{2}\right)^{n-1}=2^{3-n}$，$S_n=\frac{4\left[1-\left(\frac{1}{2}\right)^n\right]}{1-\frac{1}{2}}=8-2^{3-n}$，则 $a_n+S_n=2^{3-n}+8-2^{3-n}=8$，故D正确.
\end{solution}
\end{question}

\begin{question}[points=6]
已知 $f(x)$ 是定义在 $\mathbb{R}$ 上的奇函数，且当 $x>0$ 时，$f(x)=(x^2-3)e^x+2$，则（ ）
\begin{choices}
  \item $f(0)=0$
  \item 当 $x<0$ 时，$f(x)=-(x^2-3)e^{-x}-2$
  \item $f(x)\ge 2$ 当且仅当 $x\ge\sqrt{3}$
  \item $x=-1$ 是 $f(x)$ 的极大值点
\end{choices}
\begin{solution}
\textbf{答案：}$ABD$

\textbf{分析：}对A，根据奇函数特点即可判断；对B，利用 $f(x)=-f(-x)$ 代入求解即可；对C，举反例 $f(-1)>2$ 即可；对D，直接求导，根据极大值点判定方法即可判断.

\textbf{详解：}对A，因为 $f(x)$ 是定义在 $\mathbb{R}$ 上的奇函数，则 $f(0)=0$，故A正确；对B，当 $x<0$ 时，$-x>0$，则 $f(x)=-f(-x)=-[(-x)^2-3]e^{-x}-2=-(x^2-3)e^{-x}-2$，故B正确；对C，$f(-1)=-(1-3)e^{-1}-2=2e^{-1}-2<0$，故C错误；对D，当 $x<0$ 时，$f(x)=-(x^2-3)e^{-x}-2$，则 $f'(x)=-(2x)e^{-x}+(x^2-3)e^{-x}=(x^2-2x-3)e^{-x}$，令 $f'(x)=0$，解得 $x=-1$ 或 $x=3$（舍去），当 $x\in(-\infty,-1)$ 时，$f'(x)>0$，$f(x)$ 单调递增；当 $x\in(-1,0)$ 时，$f'(x)<0$，$f(x)$ 单调递减，则 $x=-1$ 是 $f(x)$ 的极大值点，故D正确.
\end{solution}
\end{question}

\begin{question}[points=6]
双曲线 $C:\frac{x^2}{a^2}-\frac{y^2}{b^2}=1(a>0,b>0)$ 的左、右焦点分别是 $F_1$、$F_2$，左、右顶点分别为 $A_1$，$A_2$，以 $F_1F_2$ 为直径的圆与 $C$ 的一条渐近线交于 $M$、$N$ 两点，且 $\angle NA_1M=\frac{5\pi}{6}$，则（ ）
\begin{choices}
  \item $\angle A_1MA_2=\frac{\pi}{6}$
  \item $MA_1=2MA_2$
  \item $C$ 的离心率为 $\sqrt{13}$
  \item 当 $a=\sqrt{2}$ 时，四边形 $NA_1MA_2$ 的面积为 $8\sqrt{3}$
\end{choices}
\begin{solution}
\textbf{答案：}$ACD$

\textbf{分析：}由平行四边形的性质判断A；由 $F_1M\perp F_2M$ 且 $MO=c$ 结合 $M$ 在渐近线上可求 $M$ 的坐标，从而可判断B的正误；由中线向量结合B的结果可得 $c^2=13a^2$，计算后可判断C的正误；结合BC的结果求出面积后可判断D的正误.

\textbf{详解：}不妨设渐近线为 $y=\frac{b}{a}x$，$M$ 在第一象限，$N$ 在第三象限. 对于A，由双曲线的对称性可得 $A_1MA_2N$ 为平行四边形，故 $\angle A_1MA_2=\pi-\frac{5\pi}{6}=\frac{\pi}{6}$，故A正确；对于B，方法一：因为 $M$ 在以 $F_1F_2$ 为直径的圆上，故 $F_1M\perp F_2M$ 且 $MO=c$，设 $M(x_0,y_0)$，则 $\begin{cases}x_0^2+y_0^2=c^2\\\frac{y_0}{x_0}=\frac{b}{a}\end{cases}$，故 $\begin{cases}x_0=a\\y_0=b\end{cases}$，故 $MA_2\perp A_1A_2$，由A得 $\angle A_1MA_2=\frac{\pi}{6}$，故 $MA_2=MA_1\times\frac{\sqrt{3}}{2}$ 即 $MA_1=\frac{2\sqrt{3}}{3}MA_2$，故B错误；对于C，方法一：因为 $\overrightarrow{MO}=\frac{1}{2}(\overrightarrow{MA_1}+\overrightarrow{MA_2})$，故 $4MO^2=MA_1^2+2\overrightarrow{MA_1}\cdot\overrightarrow{MA_2}+MA_2^2$，由B可知 $MA_2=b$，$MA_1=\frac{2\sqrt{3}}{3}b$，故 $4c^2=b^2+\frac{4}{3}b^2+2\times b\times\frac{2\sqrt{3}}{3}b\times\frac{\sqrt{3}}{2}=\frac{13}{3}b^2=\frac{13}{3}(c^2-a^2)$ 即 $c^2=13a^2$，故离心率 $e=\sqrt{13}$，故C正确；对于D，当 $a=\sqrt{2}$ 时，由C可知 $e=\sqrt{13}$，故 $c=\sqrt{26}$，故 $b=2\sqrt{6}$，故四边形 $NA_1MA_2$ 的面积为 $2S_{\triangle MA_1A_2}=2\times\frac{1}{2}\times2\sqrt{6}\times2\sqrt{2}=8\sqrt{3}$，故D正确.
\end{solution}
\end{question}

\section*{三、填空题}
本题共3小题，每小题5分，共15分.

\begin{question}[points=5]
已知平面向量 $\vec{a}=(x,1)$，$\vec{b}=(x-1,2x)$，若 $\vec{a}\perp(\vec{a}-\vec{b})$，则 $|\vec{a}|=$.
\fillin[{$\sqrt{2}$}]
\begin{solution}
\textbf{答案：}$\sqrt{2}$

\textbf{分析：}根据向量坐标化运算得 $\vec{a}-\vec{b}=(1,1-2x)$，再利用向量垂直的坐标表示得到方程，解出即可.

\textbf{详解：}$\vec{a}-\vec{b}=(1,1-2x)$，因为 $\vec{a}\perp(\vec{a}-\vec{b})$，则 $\vec{a}\cdot(\vec{a}-\vec{b})=0$，则 $x+1-2x=0$，解得 $x=1$. 则 $\vec{a}=(1,1)$，则 $|\vec{a}|=\sqrt{2}$.
\end{solution}
\end{question}

\begin{question}[points=5]
若 $x=2$ 是函数 $f(x)=(x-1)(x-2)(x-a)$ 的极值点，则 $f(0)=$.
\fillin[{$-4$}]
\begin{solution}
\textbf{答案：}$-4$

\textbf{分析：}由题意得 $f'(2)=0$ 即可求解 $a$，再代入即可求解.

\textbf{详解：}由题意有 $f(x)=(x-1)(x-2)(x-a)$，所以 $f'(x)=(x-a)(x-1)+(x-1)(x-2)+(x-a)(x-2)$，因为 $2$ 是函数 $f(x)$ 的极值点，所以 $f'(2)=2-a=0$，得 $a=2$. 当 $a=2$ 时，$f'(x)=2(x-2)(x-1)+(x-2)^2=(x-2)(3x-4)$，当 $x\in(-\infty,\frac{4}{3})$ 时 $f'(x)>0$，$f(x)$ 单调递增；当 $x\in(\frac{4}{3},2)$ 时 $f'(x)<0$，$f(x)$ 单调递减；当 $x\in(2,+\infty)$ 时 $f'(x)>0$，$f(x)$ 单调递增，所以 $x=2$ 是函数 $f(x)=(x-1)(x-2)(x-a)$ 的极小值点，符合题意；所以 $f(0)=(-1)\times(-2)\times(-a)=-2a=-4$.
\end{solution}
\end{question}

\begin{question}[points=5]
一个底面半径为 $4\mathrm{cm}$，高为 $9\mathrm{cm}$ 的封闭圆柱形容器（容器壁厚度忽略不计）内有两个半径相等的铁球，则铁球半径的最大值为 $\mathrm{cm}$.
\fillin[{$\frac{5}{2}$}]
\begin{solution}
\textbf{答案：}$\frac{5}{2}$

\textbf{分析：}根据圆柱与球的性质以及球的体积公式可求出球的半径.

\textbf{详解：}圆柱的底面半径为 $4\mathrm{cm}$，设铁球的半径为 $r$，且 $r<4$，由圆柱与球的性质知 $AB^2=(2r)^2=(8-2r)^2+(9-2r)^2$，即 $4r^2-68r+145=(2r-5)(2r-29)=0$，因为 $r<4$，所以 $r=\frac{5}{2}$.
\end{solution}
\end{question}

\section*{四、解答题}
本题共5小题，共77分.解答应写出文字说明，证明过程或演算步骤.

\begin{question}[points=13]
已知函数 $f(x)=\cos(2x+\varphi)$ $(0\le\varphi<\pi)$，$f(0)=\frac{1}{2}$.

（1）求 $\varphi$;

（2）设函数 $g(x)=f(x)+f\left(x-\frac{\pi}{6}\right)$，求 $g(x)$ 的值域和单调区间.
\begin{solution}
\textbf{答案：}（1）$\varphi=\frac{\pi}{3}$
（2）值域为 $[-\sqrt{3},\sqrt{3}]$；单调递减区间为 $[-\frac{\pi}{12}+k\pi,\frac{5\pi}{12}+k\pi]\ (k\in\mathbb{Z})$；单调递增区间为 $[\frac{5\pi}{12}+k\pi,\frac{11\pi}{12}+k\pi]\ (k\in\mathbb{Z})$.

\textbf{分析：}（1）直接由题意得 $\cos\varphi=\frac{1}{2},(0\le\varphi<\pi)$，结合余弦函数的单调性即可得解；（2）由三角恒等变换得 $g(x)=\sqrt{3}\cos\left(2x+\frac{\pi}{6}\right)$，由此可得值域，进一步由整体代入法可得函数的单调区间.

\textbf{详解：}（1）由题意 $f(0)=\cos\varphi=\frac{1}{2}$，且 $0\le\varphi<\pi$，所以 $\varphi=\frac{\pi}{3}$.
（2）由（1）可知 $f(x)=\cos\left(2x+\frac{\pi}{3}\right)$，所以 $g(x)=f(x)+f\left(x-\frac{\pi}{6}\right)=\cos\left(2x+\frac{\pi}{3}\right)+\cos2x=\frac{1}{2}\cos2x-\frac{\sqrt{3}}{2}\sin2x+\cos2x=\frac{3}{2}\cos2x-\frac{\sqrt{3}}{2}\sin2x=\sqrt{3}\cos\left(2x+\frac{\pi}{6}\right)$，所以函数 $g(x)$ 的值域为 $[-\sqrt{3},\sqrt{3}]$. 令 $2k\pi\le2x+\frac{\pi}{6}\le\pi+2k\pi,k\in\mathbb{Z}$，解得 $-\frac{\pi}{12}+k\pi\le x\le\frac{5\pi}{12}+k\pi,k\in\mathbb{Z}$；令 $\pi+2k\pi\le2x+\frac{\pi}{6}\le2\pi+2k\pi,k\in\mathbb{Z}$，解得 $\frac{5\pi}{12}+k\pi\le x\le\frac{11\pi}{12}+k\pi,k\in\mathbb{Z}$. 所以函数 $g(x)$ 的单调递减区间为 $[-\frac{\pi}{12}+k\pi,\frac{5\pi}{12}+k\pi]\ (k\in\mathbb{Z})$，单调递增区间为 $[\frac{5\pi}{12}+k\pi,\frac{11\pi}{12}+k\pi]\ (k\in\mathbb{Z})$.
\end{solution}
\end{question}

\begin{question}[points=15]
已知椭圆 $C:\frac{x^2}{a^2}+\frac{y^2}{b^2}=1(a>b>0)$ 的离心率为 $\frac{\sqrt{2}}{2}$，长轴长为 $4$.

（1）求 $C$ 的方程；

（2）过点 $(0,-2)$ 的直线 $l$ 与 $C$ 交于 $A,B$ 两点，$O$ 为坐标原点，若 $\triangle OAB$ 的面积为 $2$，求 $|AB|$.
\begin{solution}
\textbf{答案：}（1）$\frac{x^2}{4}+\frac{y^2}{2}=1$
（2）$\sqrt{5}$

\textbf{分析：}（1）根据长轴长和离心率求出基本量后可得椭圆方程；（2）设出直线方程并联立椭圆方程后结合韦达定理用参数 $t$ 表示面积后可求 $t$ 的值，从而可求弦长.

\textbf{详解：}（1）因为长轴长为 $4$，故 $a=2$，而离心率为 $\frac{\sqrt{2}}{2}$，故 $c=\sqrt{2}$，故 $b=\sqrt{2}$，故椭圆方程为 $\frac{x^2}{4}+\frac{y^2}{2}=1$.
（2）由题设直线 $AB$ 的斜率不为 $0$，故设直线 $l:x=t(y+2)$，$A(x_1,y_1),B(x_2,y_2)$，由 $\begin{cases}x=t(y+2)\\x^2+2y^2=4\end{cases}$ 可得 $(t^2+2)y^2+4t^2y+4t^2-4=0$，故 $\Delta=16t^4-4(t^2+2)(4t^2-4)=4(8-4t^2)>0$ 即 $-\sqrt{2}<t<\sqrt{2}$，且 $y_1+y_2=-\frac{4t^2}{t^2+2}$，$y_1y_2=\frac{4t^2-4}{t^2+2}$，故 $S_{\triangle OAB}=\frac{1}{2}\times|2t|\times|y_1-y_2|=|t|\sqrt{(y_1+y_2)^2-4y_1y_2}=|t|\sqrt{\frac{16t^4}{(t^2+2)^2}-\frac{16t^2-16}{t^2+2}}=\frac{|t|\sqrt{32-16t^2}}{t^2+2}=2$，解得 $t^2=\frac{2}{3}$，即 $t=\pm\frac{\sqrt{6}}{3}$，故 $|AB|=\sqrt{1+t^2}|y_1-y_2|=\sqrt{1+\frac{2}{3}}\times\frac{\sqrt{32-16\times\frac{2}{3}}}{\frac{2}{3}+2}=\sqrt{\frac{5}{3}}\times\frac{\sqrt{\frac{64}{3}}}{\frac{8}{3}}=\frac{\sqrt{5}}{\sqrt{3}}\times\frac{8\sqrt{3}}{8}=\sqrt{5}$.
\end{solution}
\end{question}

\begin{question}[points=15]
如图，在四边形 $ABCD$ 中，$AB\parallel CD$，$\angle DAB=90^\circ$，$F$ 为 $CD$ 的中点，点 $E$ 在 $AB$ 上，$EF\parallel AD$，$AB=3AD$，$CD=2AD$，将四边形 $EFDA$ 沿 $EF$ 翻折至四边形 $EFD'A'$，使得面 $EFD'A'$ 与面 $EFCB$ 所成的二面角为 $60^\circ$.

（1）证明: $A'B\parallel$ 平面 $CD'F$；

（2）求面 $BCD'$ 与面 $EFD'A'$ 所成的二面角的正弦值.

\begin{center}
  \includegraphics[width=.38\linewidth]{question-17-1.jpg}
\end{center}
\begin{solution}
\textbf{答案：}（1）证明见解析
（2）$\frac{\sqrt{42}}{7}$

\textbf{分析：}（1）先应用线面平行判定定理得出 $A'E\parallel$ 平面 $CD'F$ 及 $EB\parallel$ 平面 $CD'F$，再应用面面平行判定定理得出平面 $A'EB\parallel$ 平面 $CD'F$，进而得出线面平行；（2）建立空间直角坐标系，利用已知条件将点 $B,C,D',E,F$ 的坐标表示出来，然后将平面 $BCD'$ 及平面 $EFD'A'$ 的法向量求出来，利用两个法向量的数量积公式可将两平面的夹角余弦值求出来，进而可求得其正弦值.

\textbf{详解：}（1）设 $AD=1$，所以 $AB=3$，$CD=2$，因为 $F$ 为 $CD$ 中点，所以 $DF=1$，因为 $EF\parallel AD$，$AB\parallel CD$，所以 $AEFD$ 是平行四边形，所以 $AE\parallel DF$，所以 $A'E\parallel D'F$，因为 $D'F\subset$ 平面 $CD'F$，$A'E\not\subset$ 平面 $CD'F$，所以 $A'E\parallel$ 平面 $CD'F$，因为 $FC\parallel EB$，$FC\subset$ 平面 $CD'F$，$EB\not\subset$ 平面 $CD'F$，所以 $EB\parallel$ 平面 $CD'F$，又 $EB\cap A'E=E$，$EB,A'E\subset$ 平面 $A'EB$，所以平面 $A'EB\parallel$ 平面 $CD'F$，又 $A'B\subset$ 平面 $A'EB$，所以 $A'B\parallel$ 平面 $CD'F$.
（2）因为 $\angle DAB=90^\circ$，所以 $AD\perp AB$，又因为 $AB\parallel FC$，$EF\parallel AD$，所以 $EF\perp FC$，以 $F$ 为原点，$FE$，$FC$ 以及垂直于平面 $BECF$ 的直线分别为 $x$，$y$，$z$ 轴，建立空间直角坐标系. 因为 $D'F\perp EF$，$CF\perp EF$，平面 $EFD'A'$ 与平面 $EFCB$ 所成二面角为 $60^\circ$，所以 $\angle D'FC=60^\circ$. 则 $B(1,2,0)$，$C(0,1,0)$，$D'\left(0,\frac{1}{2},\frac{\sqrt{3}}{2}\right)$，$E(1,0,0)$，$F(0,0,0)$. 所以 $\overrightarrow{BC}=(-1,-1,0)$，$\overrightarrow{CD'}=\left(0,-\frac{1}{2},\frac{\sqrt{3}}{2}\right)$，$\overrightarrow{FE}=(1,0,0)$，$\overrightarrow{FD'}=\left(0,\frac{1}{2},\frac{\sqrt{3}}{2}\right)$. 设平面 $BCD'$ 的法向量为 $\vec{n}=(x,y,z)$，则 $\begin{cases}\overrightarrow{BC}\cdot\vec{n}=0\\\overrightarrow{CD'}\cdot\vec{n}=0\end{cases}$，即 $\begin{cases}-x-y=0\\-\frac{1}{2}y+\frac{\sqrt{3}}{2}z=0\end{cases}$，令 $y=\sqrt{3}$，则 $z=1$，$x=-\sqrt{3}$，则 $\vec{n}=(-\sqrt{3},\sqrt{3},1)$. 设平面 $EFD'A'$ 的法向量为 $\vec{m}=(x_1,y_1,z_1)$，则 $\begin{cases}\overrightarrow{FE}\cdot\vec{m}=0\\\overrightarrow{FD'}\cdot\vec{m}=0\end{cases}$，即 $\begin{cases}x_1=0\\\frac{1}{2}y_1+\frac{\sqrt{3}}{2}z_1=0\end{cases}$，令 $y_1=\sqrt{3}$，则 $z_1=-1$，$x_1=0$，所以 $\vec{m}=(0,\sqrt{3},-1)$. 所以 $\cos\langle\vec{m},\vec{n}\rangle=\frac{\vec{m}\cdot\vec{n}}{|\vec{m}||\vec{n}|}=\frac{0+3-1}{\sqrt{3+3+1}\times\sqrt{1+3}}=\frac{1}{\sqrt{7}}$，所以平面 $BCD'$ 与平面 $EFD'A'$ 夹角的正弦值为 $\sqrt{1-\left(\frac{1}{\sqrt{7}}\right)^2}=\frac{\sqrt{42}}{7}$.
\end{solution}
\end{question}

\begin{question}[points=17]
已知函数 $f(x)=\ln(1+x)-x+\frac{1}{2}x^2-kx^3$，其中 $0<k<\frac{1}{3}$.

（1）证明：$f(x)$ 在区间 $(0,+\infty)$ 存在唯一的极值点和唯一的零点；

（2）设 $x_1$，$x_2$ 分别为 $f(x)$ 在区间 $(0,+\infty)$ 的极值点和零点.

（i）设函数 $g(t)=f(x_1+t)-f(x_1-t)$，证明：$g(t)$ 在区间 $(0,x_1)$ 单调递减；

（ii）比较 $2x_1$ 与 $x_2$ 的大小，并证明你的结论.
\begin{solution}
\textbf{答案：}（1）证明见解析
（2）（i）证明见解析；（ii）$2x_1>x_2$，证明见解析

\textbf{分析：}（1）先由题意求得 $f'(x)=x^2\left(\frac{1}{1+x}-3k\right)$，接着构造函数 $h(x)=\frac{1}{1+x}-3k,x>0$，利用导数工具研究函数 $h(x)$ 的单调性和函数值情况，从而得到函数的单调性，进而得证函数 $f(x)$ 在区间 $(0,+\infty)$ 上存在唯一极值点；再结合 $f(0)=0$ 和 $x\to+\infty$ 时 $f(x)$ 的正负情况即可得证 $f(x)$ 在区间 $(0,+\infty)$ 上存在唯一零点；（2）（i）由（1）$x_1+1=\frac{1}{3k}$ 和 $g'(t)=f'(x_1+t)-f'(x_1-t)$ 结合（1）中所得导函数 $f'(x_1)$ 计算得到 $g'(t)=-t\left[\frac{(x_1+t)^2}{(x_1+t+1)(x_1+1)}+\frac{(x_1-t)^2}{(x_1-t+1)(x_1+1)}\right]$，再结合 $t\in(0,x_1)$ 得 $g'(t)<0$ 即可得证；（ii）由函数 $g(t)$ 在区间 $(0,x_1)$ 上单调递减得到 $0>f(2x_1)$，再结合 $f(x_2)=0$ 和函数 $f(x)$ 的单调性以及函数值的情况即可得证.

\textbf{详解：}（1）由题得 $f'(x)=\frac{1}{1+x}-1+x-3kx^2=\frac{x^2}{1+x}-3kx^2=x^2\left(\frac{1}{1+x}-3k\right)$，因为 $x\in(0,+\infty)$，所以 $x^2>0$，设 $h(x)=\frac{1}{1+x}-3k,x>0$，则 $h'(x)=-\frac{1}{(1+x)^2}<0$ 在 $(0,+\infty)$ 上恒成立，所以 $h(x)$ 在 $(0,+\infty)$ 上单调递减，$h(0)=1-3k>0$，令 $h(x_1)=0\Rightarrow x_1=\frac{1}{3k}-1$，所以当 $x\in(0,x_1)$ 时，$h(x)>0$，则 $f'(x)>0$；当 $x\in(x_1,+\infty)$ 时，$h(x)<0$，则 $f'(x)<0$，所以 $f(x)$ 在 $(0,x_1)$ 上单调递增，在 $(x_1,+\infty)$ 上单调递减，所以 $f(x)$ 在 $(0,+\infty)$ 上存在唯一极值点. 对函数 $y=\ln(1+x)-x$ 有 $y'=\frac{1}{1+x}-1=-\frac{x}{1+x}<0$ 在 $(0,+\infty)$ 上恒成立，所以 $y=\ln(1+x)-x$ 在 $(0,+\infty)$ 上单调递减，所以 $y=\ln(1+x)-x< y|_{x=0}=0$ 在 $(0,+\infty)$ 上恒成立，又因为 $f(0)=0$，$x\to+\infty$ 时 $\frac{1}{2}x^2-kx^3=\frac{1}{2}x^2(1-2kx)<0$，所以 $x\to+\infty$ 时 $f(x)<0$，所以存在唯一 $x_2\in(0,+\infty)$ 使得 $f(x_2)=0$，即 $f(x)$ 在 $(0,+\infty)$ 上存在唯一零点.
（2）（i）由（1）知 $x_1=\frac{1}{3k}-1$，则 $x_1+1=\frac{1}{3k}$，$f'(x)=x^2\left(\frac{1}{1+x}-3k\right)$，则 $g'(t)=f'(x_1+t)-f'(x_1-t)=(x_1+t)^2\left(\frac{1}{x_1+t+1}-3k\right)-(x_1-t)^2\left(\frac{1}{x_1-t+1}-3k\right)=(x_1+t)^2\left(\frac{1}{x_1+t+1}-\frac{1}{x_1+1}\right)-(x_1-t)^2\left(\frac{1}{x_1-t+1}-\frac{1}{x_1+1}\right)=\frac{-t(x_1+t)^2}{(x_1+t+1)(x_1+1)}-\frac{t(x_1-t)^2}{(x_1-t+1)(x_1+1)}=-t\left[\frac{(x_1+t)^2}{(x_1+t+1)(x_1+1)}+\frac{(x_1-t)^2}{(x_1-t+1)(x_1+1)}\right]$，因为 $t\in(0,x_1)$，所以 $x_1-t+1>0$，所以 $\frac{(x_1+t)^2}{(x_1+t+1)(x_1+1)}+\frac{(x_1-t)^2}{(x_1-t+1)(x_1+1)}>0$，所以 $g'(t)<0$，所以函数 $g(t)$ 在区间 $(0,x_1)$ 上单调递减.
（ii）$2x_1>x_2$，证明如下：由（i）知函数 $g(t)$ 在区间 $(0,x_1)$ 上单调递减，所以 $g(0)>g(x_1)$ 即 $0>f(2x_1)$，又 $f(x_2)=0$，由（1）可知 $f(x)$ 在 $(x_1,+\infty)$ 上单调递减，$x_2\in(x_1,+\infty)$，且对任意 $x\in(0,x_2)$，$f(x)>0$，所以 $2x_1>x_2$.
\end{solution}
\end{question}

\begin{question}[points=17]
甲、乙两人进行乒乓球练习，每个球胜者得 $1$ 分，负者得 $0$ 分．设每个球甲胜的概率为 $p\left(\frac{1}{2}<p<1\right)$，乙胜的概率为 $q$，$p+q=1$，且各球的胜负相互独立，对正整数 $k\ge 2$，记 $p_k$ 为打完 $k$ 个球后甲比乙至少多得 $2$ 分的概率，$q_k$ 为打完 $k$ 个球后乙比甲至少多得 $2$ 分的概率．

（1）求 $p_3$，$p_4$（用 $p$ 表示）．

（2）若 $\frac{p_4-p_3}{q_4-q_3}=4$，求 $p$．

（3）证明：对任意正整数 $m$，$p_{2m+1}-q_{2m+1}<p_{2m}-q_{2m}<p_{2m+2}-q_{2m+2}$.
\begin{solution}
\textbf{答案：}（1）$p_3=p^3$，$p_4=p^3(4-3p)$
（2）$p=\frac{2}{3}$
（3）证明见解析

\textbf{分析：}（1）直接由二项分布概率计算公式即可求解；（2）由题意 $q_3=q^3$，$q_4=q^3(4-3q)$，联立 $\frac{p_4-p_3}{q_4-q_3}=4$，$p+q=1$ 即可求解；（3）首先 $p_{2m}-p_{2m+1}=C_{2m}^{m-1}p^{m+1}q^m$，$p_{2m+2}-p_{2m+1}=C_{2m+1}^{m}p^{m+2}q^m$，同理有 $q_{2m}-q_{2m+1}=C_{2m}^{m-1}q^{m+1}p^m$，$q_{2m+2}-q_{2m+1}=C_{2m+1}^{m}q^{m+2}p^m$，作差有 $p_{2m+1}-q_{2m+1}<p_{2m}-q_{2m}$，另一方面 $p_{2m+2}-p_{2m}=\frac{(2m+1)!}{m!(m+1)!}p^mq^m\cdot p\left(p-\frac{m}{2m+1}\right)$，且同理有 $q_{2m+2}-q_{2m}=\frac{(2m+1)!}{m!(m+1)!}p^mq^m\cdot q\left(q-\frac{m}{2m+1}\right)$，作差能得到 $p_{2m}-q_{2m}<p_{2m+2}-q_{2m+2}$，由此即可得证.

\textbf{详解：}（1）$p_3$ 为打完 $3$ 个球后甲比乙至少多得两分的概率，故只能甲胜三场，故 $p_3=C_3^3(1-p)^0p^3=p^3$；$p_4$ 为打完 $4$ 个球后甲比乙至少多得两分的概率，故甲胜三场或四场，故 $p_4=C_4^3(1-p)^1p^3+C_4^4(1-p)^0p^4=4p^3(1-p)+p^4=p^3(4-3p)$.
（2）由（1）得 $p_3=p^3$，$p_4=p^3(4-3p)$，同理 $q_3=q^3$，$q_4=q^3(4-3q)$，若 $\frac{p_4-p_3}{q_4-q_3}=4$，$p+q=1$，则 $\frac{p^3(4-3p)-p^3}{q^3(4-3q)-q^3}=\frac{3p^3(1-p)}{3q^3(1-q)}=\frac{p^3q}{q^3p}=\left(\frac{p}{q}\right)^2=4$，由于 $0<p,q<1$，所以 $p=2q=2(1-p)$，解得 $p=\frac{2}{3}$.
（3）我们有 $p_{2m}-p_{2m+1}=C_{2m}^{m-1}p^{m+1}q^m$，$p_{2m+2}-p_{2m+1}=C_{2m+1}^{m}p^{m+2}q^m$，同理有 $q_{2m}-q_{2m+1}=C_{2m}^{m-1}q^{m+1}p^m$，$q_{2m+2}-q_{2m+1}=C_{2m+1}^{m}q^{m+2}p^m$. 故 $p_{2m}-p_{2m+1}=C_{2m}^{m-1}p^{m+1}q^m\cdot 1$，$q_{2m}-q_{2m+1}=C_{2m}^{m-1}q^{m+1}p^m\cdot 1$，因为 $p>q$，所以 $p_{2m}-p_{2m+1}>q_{2m}-q_{2m+1}$，即 $p_{2m+1}-q_{2m+1}<p_{2m}-q_{2m}$. 另一方面，$p_{2m+2}-p_{2m}=\frac{(2m+1)!}{m!(m+1)!}p^mq^m\cdot p\left(p-\frac{m}{2m+1}\right)$，$q_{2m+2}-q_{2m}=\frac{(2m+1)!}{m!(m+1)!}p^mq^m\cdot q\left(q-\frac{m}{2m+1}\right)$，结合 $p\left(p-\frac{m}{2m+1}\right)-q\left(q-\frac{m}{2m+1}\right)=(p-q)(p+q)-\frac{m}{2m+1}(p-q)=(p-q)\left(1-\frac{m}{2m+1}\right)=\frac{(m+1)(p-q)}{2m+1}>0$，就能得到 $p_{2m+2}-p_{2m}>q_{2m+2}-q_{2m}$，即 $p_{2m}-q_{2m}<p_{2m+2}-q_{2m+2}$，证毕.
\end{solution}
\end{question}

